\documentclass[a4paper, 11pt]{article} % das Papierformat zuerst
\usepackage[utf8]{inputenc}
\usepackage{geometry}
\geometry{a4paper, top=25mm, left=35mm, right=25mm, bottom=25mm}
\usepackage{graphicx}
\usepackage{color}
\usepackage{epstopdf}
\usepackage[T1]{fontenc}
\usepackage{setspace}
\usepackage{tabularx}
\usepackage{blindtext}
\usepackage[ngerman]{babel} %deutsche Silbentrennung
\usepackage{titlesec}
\usepackage{enumitem}
\usepackage{ifthen}
\usepackage[colorlinks=true,linkcolor=fsi]{hyperref}

\definecolor{fsi}{RGB}{0,81,150}

\newcounter{summe}

\newcommand{\abstimmung}[4]{
	\setcounter{summe}{#3}
	\addtocounter{summe}{#4}
	\begin{flushleft}
		#1\\
	Es wird über den Antrag abgestimmt.
	\end{flushleft}
	\ifthenelse{\equal{#3}{0}\AND\equal{#4}{0}}{
	\begin{center}
		\textbf{Der Antrag wird einstimmig angenommen.}
	\end{center}
	}{
	\begin{center}
		#2 \ifthenelse{\equal{#2}{1}}{Stimme}{Stimmen} dafür, #3 \ifthenelse{\equal{#3}{1}}{Stimme}{Stimmen} dagegen, #4 \ifthenelse{\equal{#42}{1}}{Enthaltung}{Enthaltungen}\\
		\ifthenelse{#2>\value{summe}}{
		\textbf{Der Antrag ist somit angenommen.}
		}{
		\textbf{Der Antrag ist somit abgelehnt.}
		}
	\end{center}
	}
}


% VARIABLEN
\newcommand{\protokoller}{Johannes Kirchner}
\newcommand{\dateOfMeeting}{22. Juni 2015}
\newcommand{\TeXer}{Johannes Kirchner}
\newcommand{\fsiPresident}{Georg Schäfer}
\newcommand{\protocolID}{8}

\begin{document}
%deckblatt start

\doublespacing
\thispagestyle{empty}

\begin{center}
\includegraphics[width=10.0cm]{../logo_fsi.eps}

\vspace*{\fill}
{\LARGE \textbf{Protokoll der Sitzung des Fachschaftsrates \\vom \dateOfMeeting}}\\
Fachschaftsrat Informatik\\
Trier University of Applied Sciences\\
\vspace{2.5cm}
\textit{
	Protokollführer: \textbf{\protokoller} \\
	\LaTeX - Umsetzung von \TeXer\\
	am \today\\
}
Protokoll ID: \textbf{\#\protocolID}
\vfill
\end{center}

\hspace*{-35cm}
\textcolor{fsi}{\rule{64.9cm}{15pt}}
\pagebreak

\setcounter{tocdepth}{2}
\tableofcontents
\pagebreak

\section{Eröffnung}
Als Protokollführer wird \protokoller~bestimmt.\\
Der Sprecher des Fachschaftsrates \fsiPresident~eröffnet die Sitzung um 19:23 Uhr.
\\\\
\textbf{Es wird festgestellt, dass der Fachschaftsrat beschlussfähig ist.}\\
\textbf{Es fehlen: Santo Pfingsten}
\\\\
Das Protokoll der letzten Sitzung wurde verlesen und genehmigt.
%Die Genehmigung des Protokolls der letzten Sitzung wird auf die nächste Sitzung vertagt.

%\abstimmung{Es wird der Antrag gestellt, Dinge zu tun}{5}{6}{2} -- {dafür}{dagegen}{enthaltungen}

\section{Events}
Der Termin für die GameDevWeek ist bereits fesgelegt und die Räume dafür im X-Gebäude wurden reserviert. Drei Studierende nehmen an der GDW im Rahmen eines Fachseminars teil.\\
Tobias Meier und Santo Pfingsten werden sich mit Janine Kleinbauer treffen, um das Abschlussgrillen des 100-Tage-Projektes zu planen.
Fabio Gimmillaro plant, die nächste \textbf{Night of the living devs} \textit{(Notld)} noch vor der Reading-Week zu veranstalten.
Der nächste Analogspieleabend findet am 09. Juli statt, im Fachschaftsraum befinden sich von der letzten noch Spiele, den Besitzer unbekannt ist.

\section{Netzwerkinstallationveranstaltung}
Aktuell haben 51 Personen einen Platz für die NIV auf der dafür vorgesehenen Webseite reserviert.\\
Das Geld für den neuen Server der Fachschaft Informatik wurde überwiesen. Das bestelle Banner zur Bewerbung der NIV wird morgen, den 23. Juni 2015 ankommen. Es wir überlegt Rabatt zu beantragen, da erst verspätet geliefert wurde.\\
Nachdem die Anträge für Mehrfachstecker genehmigt wurden, wurden diese bereits bestellt. Weiter wurden 34 Kästen Mate bei Firma Getränke Herresthal bestellt.\\
Wie die Medienanlage der Aula bedient wird, wird beim Aufbau geklärt.
Der Fachschaftsrat der Universität Trier konnte/wollte die an ihn übermittelten Plakate nicht aufhängen.\\
Für Auf-/Abbau und Verkauf gibt es bisher zu wenige Zusagen im Doodle.

\section{Finanzen}
Philipp Dippel hat ein neues Kassenbuch eingeführt. In der Ablage befindet sich eine Liste für Anschaffungen, bald soll auch auf der Webseite ein Formular dazu eingerichtet werden.\\
Dringend benötigt werden (AA-)Batterien.\\
In den letzten Tagen ist ein leerer Kasten Mate verschwunden.\\
Zu Schrank und Sofas gibt es noch keine neuen Informationen, Dominik Petersdorf schickt ein Angebot für einen neuen Schrank an den Rat.\\
Alle Anträge, die noch dem Allgemeinen Studierendenausschuss vorlagen, sind genehmigt.\\
Marco Lochen schlägt vor, einen neuen Beamer zu kaufen. Es wird vorgeschlagen, Prof. Dr. Oechsle/den Fachbereich zu bitten, das alte Leihgerät auszutauschen.

\section{Raum}
Die Studierenden, welche eine Aufsichtszeit im Fachschaftsraum übernommen haben, haben noch keinen Raumzugang erhalten. Der Putzplan wird nicht ordnungsgemäß eingehalten, wer seinen Dienst nicht absolviert ist schwer zu kontrollieren.\\
Die alten LACK-Tische wurden durch neue ersetzt.\\
Die Pflanze wird unregelmäßig und zu selten gegossen, ist aber dennoch wohlauf.

\section{Sonstiges}
\abstimmung{Es wird der Antrag gestellt, IDs für Beschlüsse einzuführen}{12}{0}{0}
Es ergeht der Beschluss \textit{(\#FSRI\_1)}, dass alle zukünftigen Beschlüsse mit einer eindeutigen ID versehen werden.

Trotz der Unfähigkeit des Fachschaftsrates Informatik der Universität Trier Werbung für unsere Events bei ihren Studierenden zu machen, wird weiterhin Werbung für diesen an der Hochschule verteilt. Georg Schäfer setzt sich stark dafür ein. Der Gast Stefan Bodenschatz ebenso.

\abstimmung{Es wird der Antrag gestellt, keine Werbung mehr für den Fachschaftsrat Informatik der Uni an der Hochschule zu machen}{2}{6}{4}

\fsiPresident~schließt die Sitzung um 19:57 Uhr

\pagebreak
\section{Bestätigung des Protokolls}
Der Sprecher des Fachschaftsrates Informatik sowie der Protokollführer dieses Protokolls bestätigen mit Ihrer Unterschrift unter diesem Protokoll, dass selbiges inhaltlich korrekt ist und alle darin aufgeführten Beschlüsse so wie beschrieben vom Rat getragen werden.
\\

\vspace{3.5cm}
\hrulefill \hfill \hrulefill

\fsiPresident \hfill \protokoller

{\footnotesize (Sprecher)\hfill (Protokollführer)}
\end{document}
